\documentclass[12pt]{article}

%-------------------------------------------------
% Packages
%-------------------------------------------------
\usepackage{amsmath, amssymb}
\usepackage{geometry}
\usepackage[hidelinks]{hyperref}   % Clean TOC (no red boxes)
\usepackage{graphicx}

%-------------------------------------------------
% Page Setup
%-------------------------------------------------
\geometry{a4paper, margin=1in}

%-------------------------------------------------
% Title Information
%-------------------------------------------------
\title{\textbf{Lecture Scribe: CSE400 Fundamentals of Probability in Computing}\\
Lecture 8: Gaussian, Uniform, Exponential, and Gamma Random Variables}

\author{Instructor: Dhaval Patel, PhD}
\date{January 29, 2026}

%-------------------------------------------------
\begin{document}

\maketitle
\tableofcontents
\newpage

%-------------------------------------------------
\section{Introduction to Moments of Continuous Random Variables (CRV)}

To analyze the shape and characteristics of probability distributions, we utilize \textbf{central moments}, which describe how a random variable deviates from its mean.

\subsection{Definitions and Step-by-Step Derivations}

\subsubsection{The $n^{th}$ Order Central Moment}

The $n^{th}$ order central moment is defined as the expected value of the difference between the random variable $X$ and its mean $\mu_X$, raised to the power of $n$:

\[
E[(X - \mu_X)^n] = \int_{-\infty}^{\infty} (x - \mu_X)^n f_X(x)\, dx
\]

%-------------------------------------------------
\subsubsection{Case $n=0$ (Zeroth Moment)}

\textbf{Reasoning:} Any value raised to the power of zero is 1.

\[
E[(X - \mu_X)^0] = E[1] = 1
\]

%-------------------------------------------------
\subsubsection{Case $n=1$ (First Central Moment)}

\[
E[X - \mu_X] = E[X] - \mu_X
\]

Since $E[X] = \mu_X$, we obtain:

\[
\mu_X - \mu_X = 0
\]

\textbf{Conclusion:} The first central moment is always zero.

%-------------------------------------------------
\subsubsection{Case $n=2$ (Variance $\sigma_X^2$)}

\[
\sigma_X^2 = E[(X - \mu_X)^2]
\]

Expanding:

\[
\sigma_X^2 = E[x^2] - 2\mu_X E[x] + \mu_X^2
\]

Substituting $E[X] = \mu_X$:

\[
\sigma_X^2 = E[x^2] - 2\mu_X^2 + \mu_X^2
\]

Final form:

\[
\sigma_X^2 = E[x^2] - \mu_X^2
\]

%-------------------------------------------------
\subsection{Higher-Order Moments: Skewness and Kurtosis}

\subsubsection{Skewness ($n=3$)}

Skewness measures the symmetry of the PDF:

\[
C_s = \frac{E[(X - \mu_X)^3]}{\sigma_X^3}
\]

\textbf{Interpretation:}

\begin{itemize}
    \item Positive ($+$): Right-skewed distribution
    \item Negative ($-$): Left-skewed distribution
\end{itemize}

%-------------------------------------------------
\subsubsection{Kurtosis ($n=4$)}

Kurtosis measures the peakedness of the PDF:

\[
c_k = E[(X - \mu_X)^4]
\]

A large value indicates a sharp peak near the mean $\mu_X$.

%=================================================
\section{Gaussian Random Variable}

The Gaussian (Normal) random variable is a fundamental model for uncertainty and noise in computing.

%-------------------------------------------------
\subsection{Definition}

A random variable $X$ is Gaussian if its PDF is:

\[
f_X(x) = \frac{1}{\sqrt{2\pi\sigma^2}}
\exp\left(-\frac{(x-m)^2}{2\sigma^2}\right)
\]

Notation:

\[
X \sim \mathcal{N}(m, \sigma^2)
\]

where $m$ is the mean and $\sigma^2$ is the variance.

%-------------------------------------------------
\subsection{Standard Normal Distribution}

A special case:

\[
m = 0, \quad \sigma^2 = 1
\]

%-------------------------------------------------
\subsection{Standard Calculation Forms}

Gaussian integrals are complex, so standard functions are used.

%-------------------------------------------------
\subsubsection{Error Function ($erf$)}

\[
erf(x) = \frac{2}{\sqrt{\pi}}
\int_{0}^{x} \exp(-t^2)\, dt
\]

%-------------------------------------------------
\subsubsection{Complementary Error Function ($erfc$)}

\[
erfc(x) = 1 - erf(x)
\]

\[
erfc(x) = \frac{2}{\sqrt{\pi}}
\int_{x}^{\infty} \exp(-t^2)\, dt
\]

%-------------------------------------------------
\subsubsection{$\Phi$-Function (Standard Normal CDF)}

\[
\Phi(x) = \frac{1}{\sqrt{2\pi}}
\int_{-\infty}^{x}
\exp\left(-\frac{t^2}{2}\right)\, dt
\]

%-------------------------------------------------
\subsubsection{Q-Function (Gaussian Tail)}

\[
Q(x) = \frac{1}{\sqrt{2\pi}}
\int_{x}^{\infty}
\exp\left(-\frac{t^2}{2}\right)\, dt
\]

%-------------------------------------------------
\subsection{Key Identities}

\[
Q(x) = 1 - \Phi(x)
\]

\[
F_X(x) = \Phi\left(\frac{x-m}{\sigma}\right)
\]

\[
Pr(X > x) = Q\left(\frac{x-m}{\sigma}\right)
\]

\[
F_X(x) = 1 - Q\left(\frac{x-m}{\sigma}\right)
\]

%=================================================
\section{Worked Example: Gaussian Probabilities}

\textbf{Problem:}

\[
f_X(x) =
\frac{1}{\sqrt{8\pi}}
\exp\left(-\frac{(x+3)^2}{8}\right)
\]

Find probabilities in terms of Q-functions.

%-------------------------------------------------
\subsection{Step 1: Identification of Parameters}

Compare with:

\[
f_X(x)=\frac{1}{\sqrt{2\pi\sigma^2}}
\exp\left(-\frac{(x-m)^2}{2\sigma^2}\right)
\]

Thus:

\[
m = -3
\]

\[
2\sigma^2 = 8 \Rightarrow \sigma^2=4 \Rightarrow \sigma=2
\]

%-------------------------------------------------
\subsection{Step 2: Solutions}

\subsubsection{$Pr(X \le 0)$}

\[
F_X(0)=\Phi\left(\frac{0-(-3)}{2}\right)
=\Phi\left(\frac{3}{2}\right)
\]

In Q-function form:

\[
1-Q\left(\frac{3}{2}\right)
\]

%-------------------------------------------------
\subsubsection{$Pr(X>4)$}

\[
Q\left(\frac{4-(-3)}{2}\right)
=Q\left(\frac{7}{2}\right)
\]

%-------------------------------------------------
\subsubsection{$Pr(|X+3|<2)$}

\[
|X+3|<2
\Rightarrow -2<X+3<2
\Rightarrow -5<X<-1
\]

Thus:

\[
Pr(|X+3|<2)=F_X(-1)-F_X(-5)
\]

\[
=\Phi(1)-\Phi(-1)
\]

Using symmetry:

\[
Pr(|X+3|<2)=2\Phi(1)-1
\]

\[
=1-2Q(1)
\]

%=================================================
\section{Worked Example: CDF Analysis}

\textbf{Problem:}

\[
F_X(x)=x^2, \quad 0\le x\le 1
\]

and:

\[
F_X(x)=1, \quad x>1
\]

%-------------------------------------------------
\subsection{Step-by-Step Analysis}

\subsubsection{PDF}

\[
f_X(x)=\frac{d}{dx}(x^2)=2x
\]

%-------------------------------------------------
\subsubsection{Mean}

\[
\mu_X=\int_{0}^{1}x(2x)\,dx
\]

\[
=\int_{0}^{1}2x^2\,dx
=\left[\frac{2x^3}{3}\right]_0^1
=\frac{2}{3}
\]

%-------------------------------------------------
\subsubsection{Variance}

\[
\sigma_X^2=\int_{0}^{1}x^2(2x)\,dx-\mu_X^2
\]

\[
=\frac{1}{2}-\left(\frac{2}{3}\right)^2
\]

\[
=\frac{1}{18}
\]

%-------------------------------------------------
\subsubsection{Skewness}

\[
c_s=-\frac{2\sqrt{2}}{5}
\]

Left-skewed distribution.

%-------------------------------------------------
\subsubsection{Kurtosis}

\[
c_k=\frac{12}{5}
\]

Platykurtic distribution.

%=================================================
\section{Engineering Applications of Probability Models}

Uncertainty is inherent in computing and hardware.

\begin{itemize}
    \item \textbf{Network Systems:} Packet delay modeled using latency statistics.
    \item \textbf{Electronic Circuits:} Thermal noise voltage is Gaussian.
    \item \textbf{IoT and Robotics:} Kalman filters estimate true location from noisy sensor data.
    \item \textbf{Image Processing:} Denoising algorithms use Gaussian priors.
\end{itemize}

%-------------------------------------------------
\end{document}
